\section{Nietzsche for Today}

\begin{quotex}
This essay by \textbf{Julius Evola} was originally published in the journal \textit{Roma} on 28 July 1971 under the title \textit{Nietzsche ancora oggi}. This is Evola's mature statement on Nietzsche, although it is not always clear if he is describing Nietzsche or himself. Tomorrow we will post our own opinions on ``Nietzsche, 40 years later". 
\end{quotex}

In Italy it seems that an interest in \textbf{Friedrich Nietzsche} has been revived. One obvious sign of it is that Adelphi of Milan is publishing critical translations of all his works; secondly, there is the almost simultaneous appearance of two books, \emph{Nietzsche} by Adriano Romualdi, which contains a full essay on this thinker followed by a selection of passages from his writings, and then the translation from the German of an excellent systematic work, \emph{Nietzsche and the Meaning of Life} by Robert Reininger. The specific approach of the second of these books leads us to pose this question: apart from the importance that Nietzsche has as a philosopher in general, what can his ideas mean today and, more precisely, which of his ideas maintain any validity?

The relevance of this problem was brought to light by Reininger, by noting that the figure of Nietzsche also has the quality of a symbol, and that his persona embodies at the same time a cause: ``it is the cause of modern man for which one fights, of this man no longer with roots in the sacred soil of tradition, oscillating between the peaks of civilization and the abysses of barbarity, seeking himself, i.e., led to create for himself a satisfying meaning for an existence of everything pushed back to himself."

One can further specify this view concerning the problem of the man of the period of nihilism, from the ``point zero of all values", of the period in which ``God is dead", on the basis of what Nietzsche had his Zarathustra announce, and that today is notoriously translated in a common and almost banal form; of the man in the period in which all external supports fail and in which — as our philosopher already said — ``the desert grows".

In the same way, one can say that however much Nietzsche can possibly provide concerns the pure individual problem. All the formulations having a possible relation with collective and political problems are put aside, those for which many wished to see a collusion between Nietzschean doctrines and some past political movements, especially Hitlerian National Socialism and which were also accused of having fueled the pride of a presumed ``Herrenvolk" (i.e., a master race) and the fixation with a poorly understood biological racism.

If a ``superman" undoubtedly constitutes a central idea of the whole of Nietzschean thought, it is in terms of a ``positive superman", it is not that grotesqueness in the style of d'Annunzio, nor the ``blond beast of prey" (this is one of Nietzsche's poorest expressions) and not even the exceptional individual who incarnates a maximum of the ``will to power", ``beyond good and evil", however without any light and without a higher sanction.

The positive superman, which suits the ``better Nietzsche", is instead to be identified with the human type who even in a nihilistic, devastated, absurd, godless world knows how to stand on his feet, because he is capable of giving himself a law from himself, in accordance with a new higher freedom.

Here we must note the clean line of demarcation that exists between Nietzsche as ``destroyer", the smasher of idols, and ``immoralist" (this latter designation which he often claimed, but only to cause a sensation: because his disdain only concerned ``petty morality" and ``herd morality"), and that ``revolution of nothing" [i.e., 1968], that anarchism from below which the profound crisis of the modern world is bringing about. It is as significant, as it is natural, that Nietzsche is absolutely unknown by the so-called ``protest" movements of today, while he was the first and the greatest of rebels. There is no correspondence in the human subject, the true elective affinities — i.e., the plebeian — of such current movements is revealed in their frequent collusion with Marxism and its offshoots, and with every social and racial slum near to the violent and destructive surface of the purely subpersonal and naturalistic strata of being.

The words of Nietzsche's Zarathustra are current and pertinent, in this regard, when he asks who strives to loosen himself from every chain: ``You call yourself free, but that does not matter to me — I ask you: free for what?", remembering that there are cases in which the only value that he possessed was thrown away together with the chain. This is a clear warning for those who today only know how to speak of ``repressions" and who feed a hysterical intolerance for every type of authority — and they feed such an intolerance — just for this reason: because they do not have in themselves a higher principle that commands.

Now the Nietzschean type, who has put ``nihilism behind him", who, in fact, ``knows how to obtain a healthy remedy from such poison", is the one who instead possesses that principle, and who therefore also knows how to give himself a law. Reininger, in this regard, is correct in seeing in Nietzsche the affirmer of an ``absolute" morality like Kant's, and certain connections could even be established with ancient Stoic ethics, which likewise advocated an interior sovereignty.

Certainly, the multiplicity of dramatically changing positions, sometimes even contradictory, among which Nietzsche sought to find his own way, can lead toward a quite different direction: for example, when Nietzsche promotes the exaltation of ``life" or when he invokes the ``fidelity to the Earth". Fidelity also to oneself: to be and to will to be what one is, sometimes this is proposed as the only possible and valid standard in the ``desert that grows". The adequate, but dangerous, standard, known even in classical antiquity before any ``existentialism".

The fundamental problem, of the greatest importance for what the best of Nietzsche can offer today, entails this danger. After what has been said at this point, one must be one's own law to oneself, it is a question of seeing what the individual finds in himself and accepting the limit reached by the multiple processes of spiritual dissolution that have acted in recent times: to see whether one finds in oneself that natural disgust for vulgarity and for every base interest, that will for a voluntary, clear discipline, that ability to freely establish ``values" and to achieve them without giving up whatever the cost, those values that in Nietzsche define the ``Overcomer" (Überwinder), the man not broken among so many things that are broken today.



\flrightit{Posted on 2012-12-30 by Cologero }

\begin{center}* * *\end{center}

\begin{footnotesize}\begin{sffamily}



\texttt{Mihai on 2012-12-30 at 11:13 said: }

I agree with your introduction, that Evola seems to combine a description of Nietzsche with a description of his own ideas. 

In Nietzsche, there is nothing of the ``positive superman" to be found- that is the person who can give himself a higher law according to a higher function, because for Nietzsche, there is no higher law or function to begin with. For him everything begins and ends in the individual- just like in every modern philosophical movement one might add. 

If Nietzsche's writings are to be exploited, then one must keep his best ideas as regarding the means (such as his view of freedom), but totally discard the ends, because they are truly dead (ends). 

It is also to be remarked that Evola only gives ambiguous explanations and descriptions whenever speaking of ``higher values" or of ``Transcendence".


\hfill

\texttt{X on 2012-12-30 at 12:12 said: }

When Evola could address such person as Nietzsche having ``an `absolute' morality" in the last years of his life, it just can represent the actual level of his own morality. It was from such level of morality he could write those words about Christianity.


\hfill

\texttt{Avery on 2012-12-30 at 18:04 said: }

Evola does have a few descriptions of higher values. He makes a few vague lists in Ruins which I will not go hunting for, and elaborates on some of them, like this one: ``Against all forms of resentment and social competition, every person should acknowledge and love his station in life, thus acknowledging the limits within which he can develop his potential; and should give an organic sense to his life and achieve its perfection, since an artisan who perfectly fulfills his function is certainly superior to a king who does not live up to his dignity." (Ruins, IT edition, 171)

However, citing the Tao Te Ching, he insists that the laws of the traditional era ``could not be reduced in any way to the domain of morality in the current sense of the word." (Revolt, 56) This reflects some statements in Guenon as well.

The main difference between Evola and Guenon, as articulated here before, is that Evola sees the articulation of Tradition as a will to solar power, while Guenon sees it as the discovery and maintenance of alchemical marriage. Of course the historical eras they most preferred were different as well.


\hfill

\texttt{Jason-Adam on 2012-12-30 at 18:40 said: }

It will be interesting to read Cologero's thoughts on this article – he is correct to see that in discussing Nietzsche, Evola is really discussing himself, and in this way we can see the total and absolute affinity Evola felt with Nietzsche, something that began at a very young age for him, Nietzsche filled for the void for him left after he abandoned Catholicism following his reading of Merejkowski's books. Nietzsche was arguable Evolas greatest influence considering how he modified the thought he took from others, even Guenon, to fit the Nietzschean pattern. As one website puts it, Evola's philosophy is an attempt to given Nietzsche a universal and metaphsyical basis.


\hfill

\texttt{Cologero on 2012-12-30 at 20:18 said: }

Regarding the artisan, Aquinas makes a similar point. That passage from Ruins, moreover, seems to contain a lot of ``moralizing". Even the Tao Te Ching can be translated as the Way of Virtue. I guess the key is to know what the current meaning of morality is.

There cannot be incompatible ways to ``see" tradition. Those approaches must be harmonized or else one must be rejected.

Which historical eras did each prefer, assuming that makes sense? Evola was on record in saying that the way he thinks was the same as that of the well bred man before the French Revolution. Guenon is on record in saying that the restoration of Tradition in the West should use the Middle Ages as the model and that form it would take should be based on the spiritual tradition at that time. We seem to be isolated in taking that project seriously.


\hfill

\texttt{Cologero on 2012-12-30 at 20:38 said: }

Dealing with the Nietzsche piece, Jason-Adam, will not bring any pleasure. I don't think it correct to say that there is a total and absolute affinity between Evola and Nietzsche. Evola often used the phrases ``the best Nietzsche" or ``the worst Nietzsche", so he clearly became aware at some point of the incompatibility of Nietzsche with Tradition. There is really no explicit doctrine of N. that E. adopts. E.g. in Revolt, we read ``the superman notion which is Nietzsche at his worst." In this piece, E. rejects ``the blonde beast of prey". E. never mentions, AFAIK, the eternal return, E. rejects N.'s naturalism, and so on. It was more a question of style. E. adopts a pattern, as you put it, a way of looking at things, and the bombastic style with its ideals of pagan greatness and anti-Christianism. This leads to some absurdities such as his wanting to praise the Middle Ages while rejecting its underlying spirituality.

Also in Revolt, E claims N had an absolute morality without knowing it. I'm sure N would deny it, had he heard that. The problem, as you accidentally point out, is that Evola had a personal philosophy instead of an impersonal metaphysic as explained by Guenon. By letting the camel's nose of Nietzscheanism under the tent of Tradition, Evola has led many astray. Yet on balance, I still think Guenon was Evola's greatest influence; E. referred to him as the ``Master of the 20th century". On balance, then, I think the best practice is to read Evola through Guenon to get the most out of his writings.


\hfill

\texttt{Jason-Adam on 2012-12-31 at 00:04 said: }

I agree with your observations, Cologero, and stand corrected.

Another point regarding Evola should be meditated on – his statement to Abd-el-Wahid Pallavicini that he was interested in temporal power moreso than spiritual. So it is fair to say that Evola places a greater emphasis on politics whereas Guenon sticks more with principles. I wonder what Evola (or Guenon) would think of Dugin's statement that one can not be a Traditionalist without personally being involved political anti-liberal activism. 

I agree with the principle of Gornahoor on taking the Moyen Age as our base to which we should return to, as opposed

to following a foreign form or waiting for a new form to arise. The road is narrow and the way is rough however…..

Not to sound tangental, but how would you respond to Dugin's words in aninterview early this year where he said that he believes Orthodoxy is a fully valid form for the east, but as for the west it is up for westerners to decide which form is best suited for them; he personally prefering neopaganism because he fears a revitalised Catholicism would attempt to impose union with the eastern churches. We went over a lot of this before in discussing Romanides who is an important if uncredited source for alot of Dugin's beliefs – but I do see myself that the east-west schism is something that can not be undone so easily as there are cultural and geopolitical impasses in addition to the doctrinal…….I hope you can suggest something I am not aware of….


\hfill

\texttt{Mihai on 2012-12-31 at 03:29 said: }

Dugin is hardly a worthy source for anything…


\hfill

\texttt{Avery on 2012-12-31 at 11:52 said: }

Yes, both were trying to work with all pre-modern ages, but where Guenon saw Tradition as something that flowered with the intellectual glory of the Middle Ages, Evola was big on Roman pre-intellectual intuition, and disapproved of Catholic Christian developments as a step downhill that opened the path to failure. He was a well-bred and orthodox pagan, and he would have made a fine Japanese, but he never made any pretense at embracing developed, medieval traditions as Guenon clearly did.

I agree that at least one of these views is wrong, or at least needs repair. Evola's use of Nietzsche and Otto Weininger can clearly be improved on. But both writers presented internally complete systems and there are good arguments to be made for both ``Tradition as the primeval pagan imperialism" and ``Tradition as fulfilled by the medieval nations". I made a study of both views in my recent book but failed to reach a conclusion; we'll see how the editor likes it.


\hfill

\texttt{Janus on 2012-12-31 at 17:08 said: }

It's important to not fall into a ``will the real Tradition please stand up" trap. Both ancient Rome and medieval Europe incarnated the Tradition in it various aspects. The difference, as Lings pointed out, what that the Christian embodiment of Tradition pointed to Love as a path to salvation, whereas this was not emphasized as much in Rome. Given that we are working in the western context, It's useful to look at both Christendom and Rome as predecessors in Tradition. But its re-manifestation won't look exactly like Christendom, nor Rome. Once an age of Tradition is lost, it can only be reborn, not resurrected into a zombie.


\hfill

\texttt{Cologero on 2012-12-31 at 18:55 said: }

Excellent points, Janus. I think we are just interested in what features are Traditional so they can be recognized. The next cycle will begin first in consciousness, so there will first appear a vanguard with that clear inner vision.


\hfill

\texttt{Fourth Magi on 2021-12-30 at 10:23 said: }

Nietzsche is great, but logic is not his forte. Otherwise, he would have noticed that humanity rejected God when it realized that the concept of God contradicts the concepts of human good and evil. There is no need to overthrow God, since God has become superhuman and needs protection from everything that is too human.


\end{sffamily}\end{footnotesize}
